\chapter{Operations}
\label{chap:operations}

\section{Building this manual}

The manual lives in \filepath{docs/system}. Build it with:

\begin{lstlisting}[language=bash]
cd docs/system
make
\end{lstlisting}

The output is \filepath{docs/system/sysdoc.pdf}. The build uses
\texttt{pdflatex} twice so the table of contents and references settle.

If the build fails, first check that the LaTeX distribution has the standard
packages used by \filepath{sitasdoc.cls}: \texttt{geometry}, \texttt{fancyhdr},
\texttt{titlesec}, \texttt{titletoc}, \texttt{listings}, \texttt{xcolor}, and
\texttt{hyperref}. The manual avoids project-specific generated
assets so it can be rebuilt from text alone.

The documentation build does not require a runnable Sitas binary, Docker, or
generated diagrams.

\section{Cleaning generated files}

\begin{lstlisting}[language=bash]
cd docs/system
make clean
make distclean
\end{lstlisting}

\texttt{make clean} removes LaTeX auxiliary files. \texttt{make distclean} also
removes the generated PDF.

\section{Linux validation}

Some runtime behavior is Linux-specific. The Docker helper mounts the project
directory into a Linux container and can request an \texttt{io\_uring}-capable
configuration:

\begin{lstlisting}[language=bash]
SITAS_DOCKER_IO_URING=1 tools/linux-docker.sh cargo test -p sitas-core --features std uring -- --nocapture
\end{lstlisting}

When the host or container configuration cannot provide \texttt{io\_uring}, the
Linux-only examples report that the backend is unavailable rather than faking a
successful run.

The Docker path provides repeatable validation for hard CPU affinity,
\texttt{epoll}, and \texttt{io\_uring}, which cannot be exercised honestly on
macOS.

\section{Continuous integration}

The repository contains a workflow definition at
\filepath{,github/workflows/ci.yml}. GitHub only discovers workflows below
\filepath{.github/workflows/}, so the definition is not active at its current
path. Once moved to the standard directory, it is configured for pushes and
pull requests to \texttt{main} and would run:

\begin{itemize}
    \item \texttt{rustfmt}, warning-denied \texttt{clippy},
          \texttt{cargo test}, and \texttt{cargo doc --no-deps} on Ubuntu and
          macOS;
    \item a fixed list of host examples on Ubuntu.
\end{itemize}

The definition does not build this PDF or run the bare-metal target check.


\section{Documentation conventions}

The codebase uses module comments, public API documentation, and local safety
comments extensively. These are conventions rather than a complete lint gate:
the crate roots do not currently enable \texttt{missing\_docs}, so
\texttt{cargo doc} checks that documentation builds but not that every public
item is documented.

\begin{itemize}
    \item \textbf{Public API} (\texttt{///}): document contracts, ownership,
          error behavior, and platform restrictions.
    \item \textbf{Module docs} (\texttt{//!}): explain each module's role and
          safety boundaries.
    \item \textbf{Unsafe blocks} (\texttt{// SAFETY:}): precede each unsafe
          operation with the invariant that makes it sound.
    \item \textbf{FFI blocks}: use \texttt{unsafe extern "C" \{ ... \}} as
          required by Rust Edition 2024.
    \item \textbf{Test names}: follow a descriptive
          \texttt{subject\_action\_expected\_outcome} convention.
\end{itemize}

The standard commands for checking these are part of the normal validation
workflow:

\begin{lstlisting}[language=bash]
cargo fmt --check
cargo clippy --all-targets -- -D warnings
cargo test
cargo doc --no-deps
cargo check -p sitas-core -p sitas-charlotte --target aarch64-unknown-none
\end{lstlisting}

\section{Validation strategy}

Use narrow tests first, then broaden. This keeps iteration fast while
protecting runtime invariants.

\begin{lstlisting}[language=bash,caption={Typical focused-to-broad workflow}]
cargo test task_snapshot -- --nocapture
cargo test scheduling_group -- --nocapture
cargo test serve_tcp -- --nocapture
cargo fmt --check
cargo test
cargo clippy --all-targets -- -D warnings
cargo doc --no-deps
\end{lstlisting}

For changes touching Linux-only code, use the Docker helper after local macOS
validation:

\begin{lstlisting}[language=bash,caption={Linux platform pass}]
tools/linux-docker.sh
SITAS_DOCKER_IO_URING=1 tools/linux-docker.sh sh -c \
  'cargo test -p sitas-core --features std uring && cargo run -p sitas --example os_uring'
\end{lstlisting}

The order matters. Focused tests make it easier to understand the failure.
Broad tests protect against regressions in adjacent runtime paths. This is
especially important in executor code, where a small lifecycle change can
affect timers, readiness, joins, and shutdown simultaneously.

\section{Interpreting unsupported features}

Unsupported is a valid result when the platform cannot provide a feature. CPU
placement is unsupported on non-Linux platforms. \texttt{io\_uring} may be
unavailable if the kernel, container runtime, or security profile blocks ring
setup. The project reports those cases honestly instead of hiding them behind
fake success.

A runtime experiment should make platform boundaries visible: readiness works
on several Unix backends, hard CPU affinity is Linux-specific, and
\texttt{io\_uring} is both Linux-specific and dependent on host/container
permissions.
